\documentclass[10pt]{extarticle}
\usepackage[a4paper,margin=1.5cm]{geometry}
\usepackage[T1]{fontenc}
\usepackage{amsmath,amssymb}
\usepackage{graphicx}
\usepackage[export]{adjustbox}
\usepackage{multicol}
\usepackage{enumitem}
\usepackage{booktabs,array,tabularx}
\setlength{\columnsep}{0.7cm}
\setlength{\parindent}{0pt}
\setlength{\parskip}{3pt}
\pagestyle{empty}
\makeatletter
\newcommand{\qfigdef}[4]{\@namedef{qf@file@#1}{#2}\@namedef{qf@w@#1}{#3}\@namedef{qf@wide@#1}{#4}}
\newcommand{\qf@check}[1]{\@ifundefined{qf@file@#1}{\PackageError{score2}{Unknown figure id #1}{}}{}}
\newcommand{\qf@img}[4][]{\includegraphics[width=#2\dimexpr\@nameuse{qf@w@#4}pt\relax,max width=#3,max height=0.45\textheight,#1]{\@nameuse{qf@file@#4}}}
\newif\ifqf@full
\newif\ifqf@boxed
\newsavebox{\qf@box}
\newcommand{\qf@open}{\global\qf@fulltrue\par\global\qf@boxedtrue
  \setbox\qf@box\vbox\bgroup\hsize\textwidth\linewidth\textwidth}
\newcommand{\qf@emit}[1]{\par\egroup\global\qf@boxedfalse
  \ifdim\dimexpr\ht\qf@box+\dp\qf@box\relax<#1\textheight\noindent\usebox{\qf@box}\par\else\unvbox\qf@box\fi}
\newcommand{\qfullsplit}{\ifqf@boxed\qf@emit{0.85}\fi}
\newcommand{\qf@close}{\ifqf@boxed\qf@emit{0.45}\fi\global\qf@fullfalse}
\newcommand{\qfullbegin}{\end{multicols}\qf@open}
\newcommand{\qfullend}{\qf@close\begin{multicols}{2}}
\newcommand{\qfullstart}{\qf@open}
\newcommand{\qfulldone}{\qf@close}
\newcommand{\qfig}[2][1]{\qf@check{#2}%
  \ifqf@full\qf@blk{#1}{#2}\else
  \ifnum\@nameuse{qf@wide@#2}=1 \expandafter\@firstoftwo\else\expandafter\@secondoftwo\fi
  {\end{multicols}\begin{center}\qf@img{#1}{\textwidth}{#2}\end{center}\begin{multicols}{2}}%
  {\qf@blk{#1}{#2}}\fi}
\newcommand{\qf@blk}[2]{\par\smallskip{\centering\qf@img{#1}{\linewidth}{#2}\par}\smallskip}
\newcommand{\qfigi}[2][1]{\qf@check{#2}\qf@img[valign=c]{#1}{\linewidth}{#2}}
\makeatother
\newcommand{\papersec}[1]{\par\bigskip{\centering\Large\bfseries PAPER-#1\par}\medskip}
\newcommand{\partsec}[2]{\par\medskip\noindent{\large\bfseries PART-#1: #2}\par\smallskip}
\newcommand{\sect}[1]{\par\smallskip\noindent\textbf{#1}\par}
\newcommand{\qn}[1]{\par\smallskip\noindent\textbf{#1.}~}
\newcommand{\paraq}[2]{\par\smallskip\noindent\textbf{Paragraph for Questions #1 and #2}\par}
\newenvironment{qoptions}{\begin{itemize}[leftmargin=1.8em,itemsep=1pt,topsep=2pt,label={}]}{\end{itemize}}
\newcommand{\opt}[2]{\item[(#1)] #2}
% figure registry (generated by score2_tools.py finalize)
\qfigdef{p1_phy_s1_q01_1}{JEE Advanced (2021 Series) - Test 04 (04 Apr 2021)-img/p1_phy_s1_q01_1.png}{197.6}{0}
\qfigdef{p1_phy_s1_q01_3}{JEE Advanced (2021 Series) - Test 04 (04 Apr 2021)-img/p1_phy_s1_q01_3.png}{143.4}{0}
\qfigdef{p1_phy_s1_q03_1}{JEE Advanced (2021 Series) - Test 04 (04 Apr 2021)-img/p1_phy_s1_q03_1.png}{115.5}{0}
\qfigdef{p1_phy_s1_q05_1}{JEE Advanced (2021 Series) - Test 04 (04 Apr 2021)-img/p1_phy_s1_q05_1.png}{128.5}{0}
\qfigdef{p1_phy_s2_q08_1}{JEE Advanced (2021 Series) - Test 04 (04 Apr 2021)-img/p1_phy_s2_q08_1.png}{126.1}{0}
\qfigdef{p1_phy_s2_q09_1}{JEE Advanced (2021 Series) - Test 04 (04 Apr 2021)-img/p1_phy_s2_q09_1.png}{107.4}{0}
\qfigdef{p1_phy_s2_q11_1}{JEE Advanced (2021 Series) - Test 04 (04 Apr 2021)-img/p1_phy_s2_q11_1.png}{166.0}{0}
\qfigdef{p1_phy_s3_q02_1}{JEE Advanced (2021 Series) - Test 04 (04 Apr 2021)-img/p1_phy_s3_q02_1.png}{89.3}{0}
\qfigdef{p1_phy_s3_q06_1}{JEE Advanced (2021 Series) - Test 04 (04 Apr 2021)-img/p1_phy_s3_q06_1.png}{130.3}{0}
\qfigdef{p1_chem_s1_q02_1}{JEE Advanced (2021 Series) - Test 04 (04 Apr 2021)-img/p1_chem_s1_q02_1.png}{31.6}{0}
\qfigdef{p1_chem_s1_q02_2}{JEE Advanced (2021 Series) - Test 04 (04 Apr 2021)-img/p1_chem_s1_q02_2.png}{35.5}{0}
\qfigdef{p1_chem_s1_q03_1}{JEE Advanced (2021 Series) - Test 04 (04 Apr 2021)-img/p1_chem_s1_q03_1.png}{31.6}{0}
\qfigdef{p1_chem_s1_q03_2}{JEE Advanced (2021 Series) - Test 04 (04 Apr 2021)-img/p1_chem_s1_q03_2.png}{26.4}{0}
\qfigdef{p1_chem_s1_q06_1}{JEE Advanced (2021 Series) - Test 04 (04 Apr 2021)-img/p1_chem_s1_q06_1.png}{291.1}{1}
\qfigdef{p1_chem_s1_q06_2}{JEE Advanced (2021 Series) - Test 04 (04 Apr 2021)-img/p1_chem_s1_q06_2.png}{170.8}{0}
\qfigdef{p1_chem_s2_q12_A}{JEE Advanced (2021 Series) - Test 04 (04 Apr 2021)-img/p1_chem_s2_q12_A.png}{387.7}{1}
\qfigdef{p1_chem_s2_q12_B}{JEE Advanced (2021 Series) - Test 04 (04 Apr 2021)-img/p1_chem_s2_q12_B.png}{275.5}{1}
\qfigdef{p1_chem_s3_q06_1}{JEE Advanced (2021 Series) - Test 04 (04 Apr 2021)-img/p1_chem_s3_q06_1.png}{386.8}{1}
\qfigdef{p2_phy_s1_q03_1}{JEE Advanced (2021 Series) - Test 04 (04 Apr 2021)-img/p2_phy_s1_q03_1.png}{360.9}{1}
\qfigdef{p2_phy_s1_q04_1}{JEE Advanced (2021 Series) - Test 04 (04 Apr 2021)-img/p2_phy_s1_q04_1.png}{99.7}{0}
\qfigdef{p2_phy_s1_q06_1}{JEE Advanced (2021 Series) - Test 04 (04 Apr 2021)-img/p2_phy_s1_q06_1.png}{108.0}{0}
\qfigdef{p2_phy_s2_q02_1}{JEE Advanced (2021 Series) - Test 04 (04 Apr 2021)-img/p2_phy_s2_q02_1.png}{84.0}{0}
\qfigdef{p2_phy_s2_q03_1}{JEE Advanced (2021 Series) - Test 04 (04 Apr 2021)-img/p2_phy_s2_q03_1.png}{101.0}{0}
\qfigdef{p2_phy_s2_q04_1}{JEE Advanced (2021 Series) - Test 04 (04 Apr 2021)-img/p2_phy_s2_q04_1.png}{146.0}{0}
\qfigdef{p2_phy_s2_q05_1}{JEE Advanced (2021 Series) - Test 04 (04 Apr 2021)-img/p2_phy_s2_q05_1.png}{94.0}{0}
\qfigdef{p2_phy_s2_q06_1}{JEE Advanced (2021 Series) - Test 04 (04 Apr 2021)-img/p2_phy_s2_q06_1.png}{68.3}{0}
\qfigdef{p2_phy_s3_q01_1}{JEE Advanced (2021 Series) - Test 04 (04 Apr 2021)-img/p2_phy_s3_q01_1.png}{103.5}{0}
\qfigdef{p2_phy_s3_q04_1}{JEE Advanced (2021 Series) - Test 04 (04 Apr 2021)-img/p2_phy_s3_q04_1.png}{245.2}{0}
\qfigdef{p2_phy_s3_q06_1}{JEE Advanced (2021 Series) - Test 04 (04 Apr 2021)-img/p2_phy_s3_q06_1.png}{70.7}{0}
\qfigdef{p2_chem_s1_q03_1}{JEE Advanced (2021 Series) - Test 04 (04 Apr 2021)-img/p2_chem_s1_q03_1.png}{77.3}{0}
\qfigdef{p2_chem_s1_q03_B}{JEE Advanced (2021 Series) - Test 04 (04 Apr 2021)-img/p2_chem_s1_q03_B.png}{88.5}{0}
\qfigdef{p2_chem_s1_q03_C}{JEE Advanced (2021 Series) - Test 04 (04 Apr 2021)-img/p2_chem_s1_q03_C.png}{123.1}{0}
\qfigdef{p2_chem_s1_q03_D}{JEE Advanced (2021 Series) - Test 04 (04 Apr 2021)-img/p2_chem_s1_q03_D.png}{73.7}{0}
\qfigdef{p2_chem_s1_q06_1}{JEE Advanced (2021 Series) - Test 04 (04 Apr 2021)-img/p2_chem_s1_q06_1.png}{126.7}{0}
\qfigdef{p2_chem_s1_q06_A}{JEE Advanced (2021 Series) - Test 04 (04 Apr 2021)-img/p2_chem_s1_q06_A.png}{46.5}{0}
\qfigdef{p2_chem_s1_q06_B}{JEE Advanced (2021 Series) - Test 04 (04 Apr 2021)-img/p2_chem_s1_q06_B.png}{56.5}{0}
\qfigdef{p2_chem_s1_q06_C}{JEE Advanced (2021 Series) - Test 04 (04 Apr 2021)-img/p2_chem_s1_q06_C.png}{55.6}{0}
\qfigdef{p2_chem_s1_q06_D}{JEE Advanced (2021 Series) - Test 04 (04 Apr 2021)-img/p2_chem_s1_q06_D.png}{50.4}{0}
\qfigdef{p2_chem_s2_q02_1}{JEE Advanced (2021 Series) - Test 04 (04 Apr 2021)-img/p2_chem_s2_q02_1.png}{37.3}{0}
\qfigdef{p2_chem_s2_q02_2}{JEE Advanced (2021 Series) - Test 04 (04 Apr 2021)-img/p2_chem_s2_q02_2.png}{44.1}{0}

\begin{document}
\begin{center}{\Large\bfseries JEE Advanced (2021 Series) --- Test 04 (04 Apr 2021)}\end{center}
\begin{multicols}{2}
\papersec{1}
\partsec{1}{PHYSICS}
\sect{Section-I(i)}
\qn{1} A particle of mass $m$ is moving along the $x$-axis with speed $v$ when it collides with a particle of mass $2m$ initially at rest. After the collision, the first particle has come to rest and the second particle has split into two equal-mass pieces that are shown in the figure. Which of the following statements correctly describes the speeds of the two pieces ? $(\theta>0)$
\qfig{p1_phy_s1_q01_1}
\begin{qoptions}
\opt{A}{Each piece moves with speed $v$.}
\opt{B}{Each piece moves with speed $v/2$.}
\opt{C}{One of the pieces moves with speed $v/2$, the other moves with speed greater than $v/2$}
\opt{D}{Each piece moves with speed greater than $v/2$.}
\end{qoptions}
\qn{2} Two particles having positive charges $+Q$ and $+2Q$ are fixed at equal distance $x$ from centre of a conducting sphere having zero net charge and radius $r$ as shown. Initially the switch S is open. After the switch S is closed, the net charge flowing out of sphere is
\qfig{p1_phy_s1_q01_3}
\begin{qoptions}
\opt{A}{$\dfrac{Qr}{x}$}
\opt{B}{$-\dfrac{Qr}{x}$}
\opt{C}{$\dfrac{3Qr}{x}$}
\opt{D}{$-\dfrac{3Qr}{x}$}
\end{qoptions}
\qn{3} A time varying uniform magnetic field, varying at constant rate $1$ T/sec exists in a circular region of radius $200$ m centered at $(0,-100)$. A conducting wire is placed along $y=\sin kx$, where $k=1$ rad/m, from $x=-\pi$ to $+\pi$. Find the magnitude of e.m.f. generated in the wire.
\qfig{p1_phy_s1_q03_1}
\begin{qoptions}
\opt{A}{$157$ V}
\opt{B}{$314$ V}
\opt{C}{$628$ V}
\opt{D}{Zero}
\end{qoptions}
\qn{4} A stick of uniform cross-sectional area, length $3$ m and relative density $0.4$ is hinged, to one of its ends at a point $1$ m deep inside water. When a weight of $600$ g is suspended at the other end, it rests in equilibrium making an angle of $60^\circ$ with the vertical. Determine the cross-section area of the stick.
\begin{qoptions}
\opt{A}{$90$ cm$^2$}
\opt{B}{$180$ cm$^2$}
\opt{C}{$45$ cm$^2$}
\opt{D}{$22.5$ cm$^2$}
\end{qoptions}
\qn{5} A prism of angle $A$ and refractive index $2$ is surrounded by medium of refractive index $\sqrt3$. A ray is incident on side PQ at an angle of incidence $i$ $(0\le i\le90^\circ)$ as shown. The refracted ray is then incident on side PR of prism. The minimum angle $A$ of prism for which ray incident on side PQ does not emerge out of prism from side PR (for any value of $i$) is
\qfig{p1_phy_s1_q05_1}
\begin{qoptions}
\opt{A}{$30^\circ$}
\opt{B}{$45^\circ$}
\opt{C}{$60^\circ$}
\opt{D}{$120^\circ$}
\end{qoptions}
\qn{6} Find the temperature of equilibrium of a perfectly black disc exposed normally to the Sun's ray on the surface of Earth. Imagine that it has a nonconducting backing so that it can radiate only to hemisphere of space. Assume temperature of surface of Sun $=6200$ K, radius of sun $=6.9\times10^8$ m, distance between the Sun and the Earth $=1.5\times10^{11}$ m. Stefan's constant $=5.7\times10^{-8}$ W/m$^2$.K$^4$
\begin{qoptions}
\opt{A}{$420.5$ K}
\opt{B}{$430$ k}
\opt{C}{$210$ k}
\opt{D}{$105$ k}
\end{qoptions}
\sect{Section-I(ii)}
\qn{7} If the potential difference applied across a Coolidge tube is decreased,
\begin{qoptions}
\opt{A}{the wavelength of the K$_\alpha$ line will remain same}
\opt{B}{the wavelength of the K$_\beta$ line will increase}
\opt{C}{the difference in wavelength between the K$_\alpha$ and K$_\beta$ lines will remain same}
\opt{D}{none of the above}
\end{qoptions}
\qn{8} Coefficient of friction between the two blocks is $0.3$ whereas the surface $AB$ is smooth. $(g=10\ \text{ms}^{-2})$
\qfig{p1_phy_s2_q08_1}
\begin{qoptions}
\opt{A}{Acceleration of the system of masses is $5.866$ ms$^{-2}$}
\opt{B}{Tension $T_1$ in the string is $17.7$ N}
\opt{C}{Tension $T_2$ in the string is about $41.4$ N}
\opt{D}{Acceleration of $10$ kg mass is $7.55$ ms$^{-2}$}
\end{qoptions}
\qn{9} In the figure shown, the plank is being pulled to the right with a constant speed $v$. If the cylinder does not slip then :
\qfig{p1_phy_s2_q09_1}
\begin{qoptions}
\opt{A}{the speed of the centre of mass of the cylinder is $2v$.}
\opt{B}{the speed of the centre of mass of the cylinder is zero.}
\opt{C}{the angular velocity of the cylinder is $v/R$.}
\opt{D}{the angular velocity of the cylinder is zero.}
\end{qoptions}
\qn{10} A cylindrical block of density $\rho$ is partially immersed in a liquid of density $3\rho$. The plane surface of the block remains parallel to the surface of the liquid. The height of the block is $60$cm. The block performs SHM when displaced from its mean position. [Use $g=9.8$ m/s$^2$]
\begin{qoptions}
\opt{A}{The maximum amplitude is $20$ cm}
\opt{B}{The maximum amplitude is $40$ cm}
\opt{C}{The time period will be $2\pi/7$ seconds}
\opt{D}{None}
\end{qoptions}
\qn{11} For a certain transverse standing wave on a long string, an antinode is formed at $x=0$ and next to it, a node is formed at $x=0.10$ m. The displacement $y(t)$ of the string particle at $x=0$ is shown in figure.
\qfig{p1_phy_s2_q11_1}
\begin{qoptions}
\opt{A}{Transverse displacement of the particle at $x=0.05$m and $t=0.05$ s is $-2\sqrt2$ cm.}
\opt{B}{Transverse displacement of the particle at $x=0.04$ m and $t=0.025$ s is $-2\sqrt2$ cm.}
\opt{C}{Speed of the travelling waves that interfere to produce this standing wave is $2$ m/s.}
\opt{D}{The transverse velocity of the string particle at $x=\dfrac1{15}$ m and $t=0.1$ s is $20\,\pi$ cm/s}
\end{qoptions}
\qn{12} A ball tied to the end of a string swings in a vertical circle under the influence of gravity
\begin{qoptions}
\opt{A}{When the string makes an angle $90^\circ$ with the vertical, the tangential acceleration is zero \& radial acceleration is somewhere between maximum and minimum}
\opt{B}{When the string makes an angle $90^\circ$ with the vertical, the magnitude of tangential acceleration is maximum \& radial acceleration is somewhere between maximum and minimum}
\opt{C}{At no place in the circular motion, tangential acceleration is equal to radial acceleration}
\opt{D}{Throughout the path whenever radial acceleration has its extreme value, the tangential acceleration is zero.}
\end{qoptions}
\sect{Section-II}
\qn{1} Sharp image of extended object which is placed perpendicular to the principle axis of a lens is $\eta$ times that of the object for a particular position of object on a screen. Without disturbing the position of object and screen, by shifting lens a position can be obtained where the sharp image is $1/\eta$ times that of object. Ratio of difference between the two positions of lens to the focal length of lens for $\eta=5$ is:
\qn{2} In the figure shown, segments AB, BC, CD and DA, as well as the diagonal segments AC and BD have a resistance $2\Omega$ each. The mid points of the diagonal wires are soldered with each other. Find the effective resistance between the mid points of sides AD and CD.
\qfig{p1_phy_s3_q02_1}
\qn{3} A capacitor of capacitance $10\,\mu$F is charged up to a potential difference of $2$ V and then the cell is removed. Now it is connected to a cell of emf $4$ V and is charged fully. In both cases the polarities of the two cells are in the same directions. Total heat produced in the charging process by $4$V cell is : (in mJ)
\qn{4} In order to impart an angular velocity to an earth satellite, the geomagnetic field can be used. Find the maximum possible angular velocity (in $10^{-2}$ rad/s) about its own axis gained by the satellite if a storage battery with a capacity of $Q=5$ Amp--hours is discharged suddenly through a coil of $N=20$ turns wound around the satellite's surface along the circumference of the largest circle. The satellite has a mass of $m=10^3$ kg and is a thin walled uniform sphere. The geomagnetic field is parallel to the winding plane and its flux density is $B=0.5$ Gauss. (1 Gauss $=10^{-4}$ Tesla)
\qn{5} Interference fringes were produced using white light in a double slit arrangement. When a mica sheet of uniform thickness of refractive index $1.6$ (relative to air) is placed in the path of light from one of the slits, the central fringe moves through some a distance. This distance is equal to the width of $30$ interference bands if light of wavelength $4800$ \AA\ is used. The thickness (in $\mu$m) of mica is :
\qn{6} A box of dimensions $1.5$ m $\times\,0.75$ m $\times\,0.75$ m and mass $250$ kg rests on a truck with the smaller face down. If the acceleration of the truck is gradually increased what is the minimum value of friction coefficient for which the box topples before it slides ?
\qfig{p1_phy_s3_q06_1}
\partsec{2}{CHEMISTRY}
\sect{Section-I(i)}
\qn{1} The bond enthalpies of C--C, C=C and C$\equiv$C bonds are $348$, $610$ and $835$ kJ/mol, respectively, at $298$ K and $1$ bar. The enthalpy of polymerization per mole of 2-Butyne at $298$ K and $1$ bar, as shown below, is
\[n\,\mathrm{CH_3C{\equiv}CCH_3(g)}\longrightarrow\left[\mathrm{CH_2{-}CH{=}CH{-}CH_2}\right]_n\mathrm{(g)}\]
\begin{qoptions}
\opt{A}{$-123$ kJ}
\opt{B}{$-132$ kJ}
\opt{C}{$-139$ kJ}
\opt{D}{$-37$ kJ}
\end{qoptions}
\qn{2} Consider the following statements:\\
(I) Halogen have highest 1$^{\mathrm{st}}$ Ionization energy in the respective period\\
(II) Electron gain enthalpy of alkali metal has positive value\\
(III) CrO$_3$ and Mn$_2$O$_7$ are acidic oxide\\
(IV) In sp$^3$d$^2$ hybridization set of d$_{x^2-y^2}$ and d$_{z^2}$ orbitals is involved\\
Which of the above statement incorrect one ?
\begin{qoptions}
\opt{A}{II and IV}
\opt{B}{III and IV}
\opt{C}{I, II and III}
\opt{D}{I and II}
\end{qoptions}
\qn{3} \begin{center}
(I)~\qfigi{p1_chem_s1_q02_1}\quad (II)~\qfigi{p1_chem_s1_q02_2}\par\medskip
(III)~\qfigi{p1_chem_s1_q03_1}\quad (IV)~\qfigi{p1_chem_s1_q03_2}
\end{center}
The correct order of stability for the given canonical structure is
\begin{qoptions}
\opt{A}{IV $>$ II $>$ III $>$ I}
\opt{B}{IV $>$ II $>$ I $>$ III}
\opt{C}{II $>$ IV $>$ III $>$ I}
\opt{D}{II $>$ IV $>$ I $>$ III}
\end{qoptions}
\qn{4} For the reaction $2\mathrm{NO(g)}+\mathrm{H_2(g)}\longrightarrow\mathrm{N_2O(g)}+\mathrm{H_2O(l)}$\\
\textbf{Experiment 1:} The value of $-\dfrac{dp}{dt}$ was, obtained by experiment(s), $2.25$ torr/sec for a pressure of $450$ torr of NO(g) and $0.25$ torr/sec for a pressure of $150$ torr, the pressure of H$_2$(g) being constant.\\
\textbf{Experiment 2:} The value of $-dp/dt$ was $1.90$ torr/sec for H$_2$(g) pressure of $291$ Torr and $0.951$ torr/sec for H$_2$(g) pressure of $145$ torr the pressure of NO(g) being constant.\\
Which analytical statement is correct on the basis of above experiment(s)\\
(I) Order of reaction is $3$\\
(II) Order with respect to NO(g) and H$_2$(g) are $3$ and $0$ respectively\\
(III) Order with respect to H$_2$(g) and NO(g) are $1$ and $2$ respectively\\
(IV) Rate of reaction depend upon the concentration/pressure of reactant(s)
\begin{qoptions}
\opt{A}{I, II and IV}
\opt{B}{I and III}
\opt{C}{I, III and IV}
\opt{D}{I only}
\end{qoptions}
\qn{5} Which of the following pairs correctly represent the comparision of their dipole moment?
\begin{qoptions}
\opt{A}{CH$_3$Cl $<$ CH$_3$F}
\opt{B}{CH$_3$Cl $<$ CHCl$_3$}
\opt{C}{NH$_3$ $<$ NF$_3$}
\opt{D}{HCl $<$ HF}
\end{qoptions}
\qfullbegin
\qn{6} What are the possible product(s) in are given chemical reaction.
\qfig{p1_chem_s1_q06_1}
\qfullsplit
\qfig{p1_chem_s1_q06_2}
\begin{qoptions}
\opt{A}{$1$ and $2$}
\opt{B}{$3$ and $4$}
\opt{C}{$3$, $4$, $5$ and $6$}
\opt{D}{$1$, $2$, $3$, $4$, $5$ and $6$}
\end{qoptions}
\qfullend
\sect{Section-I(ii)}
\qn{7} Which of the following is/are correct statement(s).
\begin{qoptions}
\opt{A}{Lyophilic solution are reversible whereas lyophobic solution are irreversible}
\opt{B}{Dipersed phase shows no affinity with dispersion medium in case of lyophilic colloidal}
\opt{C}{Dipersed phase molecules of lyophilic colloidals are large enough to be in colloidal range}
\opt{D}{Dispersed phase molecules of lyophobic colloidals are large enough to be close to upper limit of colloidal range}
\end{qoptions}
\qn{8} Which of the following statement(s) is/are correct?
\begin{qoptions}
\opt{A}{In [Fe(H$_2$O)$_4$Cl$_2$]Cl, the ligands has satisfied only the secondary valencies of ferric ion}
\opt{B}{In [Fe(NH$_3$)$_3$(CN)$_3$], the ligands has satisfied both primary and secondary valencies of ferric ion}
\opt{C}{In [Fe(H$_2$O)$_5$Cl]Cl, the ligands has satisfied only the secondary valencies of ferrous ion}
\opt{D}{In [Cu(NH$_3$)$_4$]SO$_4$, the ligands has satisfied only the secondary valencies of copper ion}
\end{qoptions}
\qn{9} Select the correct statement about Grignard reagent.
\begin{qoptions}
\opt{A}{3 and 4 membered cyclic ethers can not be used as a solvent in the reaction in which Grignard reagent is one of the reactant}
\opt{B}{Aryl, vinyl, allyl group can not be taken as organic part of Grignard reagent}
\opt{C}{Grignard reagent can not be prepared from alkyl or aryl fluorides}
\opt{D}{The role of ether is to solvolyse RMgX and also provide medium for reaction}
\end{qoptions}
\qn{10} Acetone and carbon disulphide form binary liquid solution showing positve deviation from Raoult's law. The normal boiling point (T$_b$) of pure acetone is less than that of pure CS$_2$. Pick out the \textbf{incorrect} statments among the following.
\begin{qoptions}
\opt{A}{Boiling temperature of mixture is always less than boiling temperature of acetone}
\opt{B}{Boiling temperature of Azeotropic mixture is always less than boiling temperature of pure CS$_2$}
\opt{C}{When a small amount CS$_2$ (less volatile component) is added to excess of acetone boiling point of resulting mixure increase}
\opt{D}{A mixture of CS$_2$ and CH$_3$COCH$_3$ can be completely separated by simple fractional distillation}
\end{qoptions}
\qn{11} Select the correct statement:
\begin{qoptions}
\opt{A}{Kaolinite is the ore of aluminium}
\opt{B}{Extraction of lead from galena involves roasting in limited supply of air at moderate temperature followed by self reduction at higher temperature (to melt the charge)}
\opt{C}{Extraction of zinc from zinc blende involves roasting followed by reduction with carbon}
\opt{D}{Major chemical component of `slag' formed during the extraction of iron and copper is FeSiO$_3$}
\end{qoptions}
\qfullbegin
\qn{12} Which of the following statement(s) and reaction(s) are correct:
\begin{qoptions}
\opt{A}{\qfigi{p1_chem_s2_q12_A}}
\opt{B}{\qfigi{p1_chem_s2_q12_B}}
\opt{C}{Lactose is a disaccharide that has one acetal and one hemiacetal unit}
\opt{D}{Reaction between carbonyl compound and monohydric alcohol is not favourable in acetal formation}
\end{qoptions}
\qfullend
\sect{Section-II}
\qn{1} $0.001$ mM, $100$ ml HCl is titrated with $0.001$ mM NaOH and conductance continuously measured, what will be the molar conductance at end point of titration.\\
Given that $\lambda^{\circ}_{\mathrm{H^+}}$, $\lambda^{\circ}_{\mathrm{Na^+}}$, $\lambda^{\circ}_{\mathrm{Cl^-}}$ and $\lambda^{\circ}_{\mathrm{OH^-}}$ are $349.8$, $50.11$, $76.34$ and $198.50$ Mho cm$^2$ mol$^{-1}$ respectively.
\qn{2} Mn$^{+2}$, Fe$^{+2}$, Ni$^{+2}$, Cu$^{+2}$ ions form different type of inner and outer orbital complexes like [Mn(SCN)$_6$]$^{4-}$, [Fe(CN)$_6$]$^{4-}$, [Ni(NH$_3$)$_6$]$^{+2}$, [Cu(NH$_3$)$_6$]$^{+2}$. What is the maximum magnetic moment from among them.
\qn{3} In the following sequence, product P, Q, R and S are formed, if the molar mass of S is M, then what is value of $\dfrac{M}{10}$.
\[\begin{gathered}\mathrm{PhCH_2C{\equiv}CCH_2CHO}\ \xrightarrow[\text{(ii) PBr}_3]{\text{(i) NaBH}_4}\ \mathrm{P}\ \xrightarrow[\substack{\text{(ii) CO}_2\\\text{(iii) H}_3\text{O}^{\oplus}}]{\text{(i) Mg/ether}}\ \mathrm{Q}\\ \xrightarrow{\text{SOCl}_2}\ \mathrm{R}\ \xrightarrow[\text{(ii) Pd/BaSO}_4\text{/quinoline}]{\text{(i) H}_2\text{(g)}}\ \mathrm{S}\end{gathered}\]
\qn{4} Calculate $\Delta G_{\text{reaction}}$ (kJ/mol) for the given reaction at $300$ K.\\
$\mathrm{A_{2(g)}+B_{2(g)}\rightleftharpoons2AB_{(g)}}$\\
and at partial pressure of $10^{-2}$ bar, $10^{-1}$ bar and $10^{-4}$ bar respectively for A$_2$, B$_2$ and AB.\\
\textbf{Given:}\\
$\Delta H_f^{\circ}$ AB $=180$ kJ/mol ; $\Delta H_f^{\circ}$ A$_2$ $=60$ kJ/mol\\
$\Delta H_f^{\circ}$ B$_2$ $=29.5$ kJ/mol ; $\Delta S_f^{\circ}$ AB $=210$ J/K mol\\
$\Delta S_f^{\circ}$ A$_2$ $=190$ J/K mol ; $\Delta S_f^{\circ}$ B$_2$ $=205$ J/K mol\\
Use: $2.303\,R\times300=5750$ J/mole
\qn{5} $4$ gm calcium burns in excess nitrogen to produce a white powder which dissolves in sufficient water to produce a gas (A) and an alkaline solution. The solution on bubbling dry air produces a solid (B) on the surface. Calculate amount of B in gm. (assume $100\%$ conversion in each reaction)
\qfullbegin
\qn{6} In the following sequence the products are P, Q, R, S and T
\qfig{p1_chem_s3_q06_1}
\qfullsplit
$\therefore$ T is a red colour compound, What is the degree of unsaturation in T.
\qfullend
\partsec{3}{MATHEMATICS}
\sect{Section-I(i)}
\qn{1} The locus of centre of the circle of radius $3$ which touches the locus of P obtained by $2\text{PA}=3\text{PB}$ \{where $\text{A}\left(9\sqrt2,9\sqrt2\right)$ and $\text{B}\left(4\sqrt2,4\sqrt2\right)$ are given fixed points\} externally, is :
\begin{qoptions}
\opt{A}{$x^2+y^2=81$}
\opt{B}{$x^2+y^2=144$}
\opt{C}{$x^2+y^2=225$}
\opt{D}{$x^2+y^2=324$}
\end{qoptions}
\qn{2} The number of integral solutions for $a,b,c$ and $d$ of the equation $a\cdot b\cdot c\cdot d=2021$ is :
\begin{qoptions}
\opt{A}{$16$}
\opt{B}{$32$}
\opt{C}{$64$}
\opt{D}{$128$}
\end{qoptions}
\qn{3} If X and Y are two non-empty sets, where $f:X\to Y$ is the function, is defined such that $f(C)=\{f(x):x\in C\}$ for $C\subseteq X$ and $f^{-1}(D)=\{x:f(x)\in D\}$ for $D\subseteq Y$ for any $A\subseteq X$ and $B\subseteq Y$, then
\begin{qoptions}
\opt{A}{$f^{-1}(f(A))=A$}
\opt{B}{$f^{-1}(f(A))=A$ only if $f(x)=Y$}
\opt{C}{$f(f^{-1}(B))=B$ only if $B\subseteq f(x)$}
\opt{D}{$f(f^{-1}(B))=B$}
\end{qoptions}
\qn{4} Let $\displaystyle\int\dfrac{x^2-x+2}{(x+1)(x-1)^3}\,dx=A\ln|x+1|+B\ln|x-1|+\dfrac{C}{x-1}+\dfrac{D}{(x-1)^2}+E$, where $E$ is an arbitrary constant then which of the following is correct ? ($A,B,C,D$ are real constants)
\begin{qoptions}
\opt{A}{$A+D=C+2B$}
\opt{B}{$A+D=C-2B$}
\opt{C}{$A+D=C+B$}
\opt{D}{$A+D=C-B$}
\end{qoptions}
\qn{5} Let $f:R\to R$ be such that $f(1)=\sqrt2$ and $f'(1)=2\sqrt2$, then the value of $\displaystyle\lim_{x\to0}\left(\dfrac{f(1+x)}{f(1)}\right)^{\frac1x}$ equals
\begin{qoptions}
\opt{A}{$1$}
\opt{B}{$e^{\frac12}$}
\opt{C}{$e^2$}
\opt{D}{$e^3$}
\end{qoptions}
\qn{6} If $\theta\in\left[0,\dfrac\pi{108}\right]$, then range of the function $f(\theta)=\dfrac{2\sin\theta}{\cos3\theta}+\dfrac{2\sin3\theta}{\cos9\theta}+\dfrac{2\sin9\theta}{\cos27\theta}+\tan\theta$ is :
\begin{qoptions}
\opt{A}{$\left[0,\sqrt2-1\right]$}
\opt{B}{$\left[0,\sqrt2+1\right]$}
\opt{C}{$[0,1]$}
\opt{D}{$[0,\infty]$}
\end{qoptions}
\sect{Section-I(ii)}
\qn{7} DG and CF are two poles of same height situated $5$m apart from each other such that D and C are the top of the poles and G and F are the feet of poles on the ground. Two observers stand on the line joining the feet of poles on opposite sides of the poles at points A and B, at a distance $23$ m from each other, such that A-G-F-B, in order, are collinear. If the distance between A and C is $25$m and B and D is $17$ m, then
\begin{qoptions}
\opt{A}{the distance between A \& D will be $15\sqrt2$m}
\opt{B}{the distance between A \& D will be $15\sqrt3$m}
\opt{C}{the distance between B \& C will be $9\sqrt3$m}
\opt{D}{the distance between B \& C will be $3\sqrt{26}$m}
\end{qoptions}
\qn{8} The value of $\dfrac12\cos^{-1}\dfrac35+\dfrac12\tan^{-1}\dfrac{12}5$ is equal to
\begin{qoptions}
\opt{A}{$\tan^{-1}\dfrac12+\tan^{-1}\dfrac23$}
\opt{B}{$\sin^{-1}\dfrac1{\sqrt5}+\cos^{-1}\dfrac3{\sqrt{13}}$}
\opt{C}{$\cos^{-1}\dfrac2{\sqrt5}+\sin^{-1}\dfrac2{\sqrt{13}}$}
\opt{D}{$\tan^{-1}\dfrac14+\tan^{-1}\dfrac29+\cot^{-1}\dfrac32$}
\end{qoptions}
\qn{9} Let $y=x^2+ax+1$ be a parabola such that no point of it lies below x axis. A tangent to this parabola at its point of intersection with the y axis also touches the circle $x^2+y^2=r^2$. Then -
\begin{qoptions}
\opt{A}{The number of integral values of `$a$' is $3$}
\opt{B}{The radius of the circle when `$a$' attains its maximum value is $\dfrac1{\sqrt5}$}
\opt{C}{The slope of the tangent when radius of the circle is maximum is $0$}
\opt{D}{The minimum area bounded by the tangent and the coordinate axes is $\dfrac14$}
\end{qoptions}
\qn{10} Let $z_1,z_2\in\mathbb C$ such that $|z_1-z_2|=4$ and $|z_1|<|z_2|$, then -
\begin{qoptions}
\opt{A}{$|z_1|+|z_2|=4\Rightarrow\arg\left(\dfrac{z_1}{z_2}\right)=\pi$}
\opt{B}{$|z_1|+|z_2|=|z_1+z_2|\Rightarrow\arg\left(\dfrac{z_1}{z_2}\right)=0$}
\opt{C}{$|z_1|-|z_2|=-4\Rightarrow\arg\left(\dfrac{z_1}{z_2}\right)=\pi$}
\opt{D}{$|z_2|-|z_1|=|z_1+z_2|\Rightarrow\arg\left(\dfrac{z_1}{z_2}\right)=0$}
\end{qoptions}
\qn{11} An instructor has a question bank consisting of $300$ easy true/false type questions, $200$ difficult true/false type questions, $500$ easy multiple choice questions and $400$ difficult multiple choice questions. If a question is selected at random from the question bank then the probability that the selected question is
\begin{qoptions}
\opt{A}{an easy question given that it's a multiple choice question is $\dfrac89$}
\opt{B}{a multiple choice question given that it's an easy question is $\dfrac58$}
\opt{C}{a difficult question given that it's a true/false type question is $\dfrac25$}
\opt{D}{an easy true-false question is $\dfrac3{14}$}
\end{qoptions}
\qn{12} Let A and B be invertible square matrices of order `$n$' where $n\ge3$, is an odd number. Then which of the following statement(s) is / are true ?
\begin{qoptions}
\opt{A}{adj (adjA) $=|A|^{n-2}A$}
\opt{B}{adj $(AB)^T=(\text{adj }A)^T(\text{adj }B)^T$}
\opt{C}{$|\text{adj(adj }A)|=|A|^{n-3}$}
\opt{D}{$A^2B^2-B^2A^2$ will be an invertible matrix if A and B are skew symmetric and symmetric matrices respectively.}
\end{qoptions}
\sect{Section-II}
\qn{1} A cone is made from a circular sheet of radius $\sqrt3$ by cutting out a sector and keeping the cut edges of the remaining piece together. Find the maximum volume attainable for the cone.
\qn{2} If $x\ln x>k(x-1)-\ln x\ \ \forall x>1$, then find the greatest possible such integral value of $k$.
\qn{3} Let $f(x)=\displaystyle\int_0^x\dfrac{2t+1}{t^2-2t+2}\,dt\ \forall x\in[-1,1]$. If the range of $f(x)$ be $[a,b]$ where $a,b\in R$ then $a+b=\ell n\left(\dfrac mn\right)+3\tan^{-1}\left(\dfrac pq\right)$ where $m,n,p,q\in N$ such that $m$ \& $n$ are coprime to each other so as $p$ \& $q$. Evaluate $m+n+p+q$.
\qn{4} Let $P(x_0,y_0)$ be a point on the curve $C:(x^2-11)(y+1)+4=0$ where $x_0,y_0\in N$. Find the area of the triangle formed by the normal drawn to the curve $C$ at P and the cordinate axes.
\qn{5} The solution of differential equation $\dfrac{dy}{dx}=3x^2y^2-\dfrac{6y}x$ such that $y(1)=\dfrac13$ is given by $y=\dfrac1{x^m\left(1+n\cdot x^p\right)}$ where $m,n,p\in N$. Find the value of $m^2+n^2+p^2$.
\qn{6} Let $\vec a,\vec b,\vec c$ and $\vec d$ be four distinct non-zero vectors such that $\left(\vec a\times\vec b\right)\times\left(\vec c\times\vec d\right)=m\vec c+n\vec d$ where $m$ \& $n$ are scalars. Find the value of $m+n$ if $\left(\vec a\times\vec b\right)\cdot\left(\vec c-\vec d\right)=-3$.
\papersec{2}
\partsec{1}{PHYSICS}
\sect{Section-I}
\qn{1} A solid sphere of mass $M$ is placed on a smooth horizontal surface. A sudden blow is given horizontally to the sphere at a height $h=4R/5$ above the centre line. If $I$ is the impulse of the blow.
\begin{qoptions}
\opt{A}{The minimum time after which the highest point will touch the ground is $\dfrac{MR\pi}{2I}$}
\opt{B}{The displacement of the centre of mass during this interval is $\dfrac{R\pi}{2}$}
\opt{C}{Angular velocity of sphere just after impulse provide is $\dfrac{I}{MR}$}
\opt{D}{Velocity of sphere just after impulse provide is $\dfrac{I}{M}$}
\end{qoptions}
\qn{2} The mechanical energy of a system is given by $E=Ax^2+Bv^2$ where $x$ is the displacement from $(x=0)$ and $v$ is the velocity of the body at $x$ then choose the correct statement(s). Assume $A$ and $B$ positive constant.
\begin{qoptions}
\opt{A}{The maximum $x$ coordinate is $\sqrt{\dfrac EA}$}
\opt{B}{Maximum velocity of the particle during motion is $2\sqrt{\dfrac EB}$}
\opt{C}{$x=0$ is the equilibrium position.}
\opt{D}{The acceleration of the system at any point is $\dfrac{-Ax}{B}$}
\end{qoptions}
\qfullbegin
\qn{3} Consider a film of thickness $L$ as shown in four different case blow. Notice the observation of film with perpendicular falling light. Mark the correct statement(s): (Take $L=\dfrac{3\lambda}{2}$, where $\lambda$ is wavelength)
\qfig{p2_phy_s1_q03_1}
\qfullsplit
\begin{qoptions}
\opt{A}{For (1) and (2), the reflection at film interfaces causes zero phase difference for two reflected rays.}
\opt{B}{For (2) and (3), the reflection at film interfaces causes a phase difference of $\pi$ for two reflected rays.}
\opt{C}{For (1), the film will appear dark, if it is observed through reflected rays from film interfaces.}
\opt{D}{For (3), the film will appear dark, if it is observed through reflected rays from film interfaces}
\end{qoptions}
\qfullend
\qn{4} A non-uniform magnetic field $\vec B=B_0\left(1+\dfrac yd\right)\hat k$ is present in region of space in between $y=0$ \& $y=d$. The lines are shown in the diagram. A particle of mass `$m$' and positive charge `$q$' is moving. Given an initial velocity $\vec v=v_0\hat i$. Find the components of velocity of the particle when it leaves the field.
\qfig{p2_phy_s1_q04_1}
\begin{qoptions}
\opt{A}{$V_y=\sqrt{V_0^2-\left(V_0-\dfrac{3qB_0}{2m}d\right)^2}$}
\opt{B}{$V_y=\sqrt{V_0^2-\left(V_0-\dfrac{qB_0}{2m}d\right)^2}$}
\opt{C}{$V_x=V_0-\dfrac{qB_0d}{2m}$}
\opt{D}{$V_x=V_0-\dfrac{3qB_0d}{2m}$}
\end{qoptions}
\qn{5} A parallel plate capacitor of capacitance $10$ mF is connected to a cell of emf $10$ Volt and fully charged. Now a dielectric slab $(k=3)$ of thickness equal to the gap between the plates, is very slowly inserted to completely filled in the gap, keeping the cell connected. During the filling process:
\begin{qoptions}
\opt{A}{the increase in charge on the capacitor is $200$ mC.}
\opt{B}{the heat produced is non-zero.}
\opt{C}{energy supplied by the cell $=$ increase in stored potential energy $+$ work done on the person who is filling the dielectric slab.}
\opt{D}{energy supplied by the cell $=$ increase in stored potential energy $+$ work done by the person who is filling the dielectric slab $+$ heat produced.}
\end{qoptions}
\qn{6} A capacitor of capacitance $C_1$ charged to a voltage $V_0$ is connected to another uncharged identical capacitor and an inductor of inductance $L$ with switch S as shown in the figure. At $t=0$, the switch S is closed. Then
\qfig{p2_phy_s1_q06_1}
\begin{qoptions}
\opt{A}{the natural frequency of oscillations is $\sqrt{\dfrac{2}{LC_1}}$}
\opt{B}{the natural frequency of oscillations is $2\sqrt{\dfrac{2}{LC_1}}$}
\opt{C}{the charge on the capacitor on the right is $\dfrac{C_1V_0}{2}\,(1+\cos\omega t)$}
\opt{D}{the charge on the capacitor on the left is $\dfrac{C_1V_0}{2}\,(1+\cos\omega t)$}
\end{qoptions}
\sect{Section-II}
\qn{1} Raindrops are falling with a velocity $10\sqrt2$ m/s making an angle of $45^0$ with the vertical. The drops appear to be falling vertically to a man running with constant velocity. The velocity of rain drops change such that the rain drops now appear to be falling vertically with $\sqrt3$ times the velocity it appeared earlier to the same person running with same velocity. The angle (in degrees) between the initial and the final velocity vectors of the raindrops with respect to the ground is?
\qn{2} S is a stationary source of sound emitting waves of frequency $440$ Hz. These waves are reflected by a stationary wall \& heard by observer O moving away with $v=10$ m/s. Find beat frequency (in Hz) heard by observer in the position shown (due to superposition between reflected \& direct waves) ($V_{\text{sound}}=330$ m/s)
\qfig{p2_phy_s2_q02_1}
\qn{3} A ring of mass $m=1$ kg can slide over a smooth vertical rod. A light string attached to the ring passing over a smooth fixed pulley at a distance of $L=0.7$ m from the rod as shown in the figure. At the other end of the string mass $M=5$ kg is attached, lying over a smooth fixed inclined plane of inclination angle $37^\circ$. The ring is held in level with the pulley and released. Determine the instantaneous power of the gravitational force acting on the ring (in watts) when the string makes an angle $(\alpha=37^\circ)$ with the horizontal. [$\sin37^\circ=0.6$]
\qfig{p2_phy_s2_q03_1}
\qn{4} An object is placed $21$ cm in front of a concave mirror of radius of curvature $10$ cm. A glass slab of thickness $3$ cm and refractive index $1.5$ is then placed close to the mirror in the space between the object and the mirror. Find the position of the final image (in cm) formed. (You may take the distances of the nearer surface of the slab from the mirror to be $1.0$ cm).
\qfig{p2_phy_s2_q04_1}
\qn{5} The circuit diagram given in the figure shows the experimental setup for the measurement of unknown resistance by using a meter bridge. The wire connected between the points P \& Q has non-uniform resistance such that resistance per unit length varies directly as the distance from the point P. Null point is obtained with the jockey J with $R_1$ and $R_2$ in the given position. On interchanging the positions $R_1$ and $R_2$ in the gaps the jockey has to be displaced through a distance $\Delta$ from the previous position along the wire to establish the null point. If the ratio of $\dfrac{R_1}{R_2}=3$, find the value of $\Delta$ (in cm). Ignore any end corrections. [Take $\sqrt3=1.7$]
\qfig{p2_phy_s2_q05_1}
\qn{6} Figure shows a vertical force $F$ that is applied tangentially to a uniform cylinder of weight $W$. The coefficient of static friction between the cylinder and all surfaces is $0.5$. The maximum force that can be applied without causing the cylinder to rotate is $\dfrac{3W}{n}$. Find `$n$'.
\qfig{p2_phy_s2_q06_1}
\sect{Section-III}
\qn{1} A slab of mass $m_2$ ($=10$ kg) is resting on a frictionless floor with a second mass $m_1$ ($=1$kg) on its top. The coefficient of friction between $m_1$ and $m_2$ is $\mu=0.1$. A horizontal force of $15$ N is applied to the lower block at $t=0$. The time required by the smaller block to move a distance of $1$ m on the bigger block is $(\sqrt n)$ sec. Find `$n$'. [Take $g=10$ m/s$^2$]
\qfig{p2_phy_s3_q01_1}
\qn{2} A hydrogen like atom is in an excited state `$2n$' can emit a maximum energy photon of $204$eV. If it makes a transition to a quantum state `$n$', a photon of energy $40.8$eV is emitted. Calculate `$n$'. Take ground state energy of hydrogen atom to be $-13.6$eV.
\qn{3} $2$ moles of a diatomic gas undergoes a thermodynamic process $\dfrac{PT^2}{V}=\text{constant}$. The molar heat capacity of the gas is $\dfrac{nR}{2}$. Find $n$.
\qn{4} The long, horizontal pair of rails shown in the figure is connected using resistance $R$. The distance between the rails is $\ell$, the electrical resistance of the rails is negligible. A conducting wire of mass $m$ and length $\ell$ can slide without friction on the pair of rails, in a vertical, homogeneous magnetic field of induction $B$. A force of magnitude $F_0$ is exerted for sufficiently long time onto the conducting wire, so that the speed of the wire becomes nearly constant. The force $F_0$ is now removed at a certain point P. The distance the conducting wire cover on rails from point P before stopping is $(316+n)$ meters. Find `$n$'. (Given : $F_0=20$N, $m=1.6$ gm, $R=0.01\Omega$, $\ell=10$ cm, $B=0.1$ T)
\qfig{p2_phy_s3_q04_1}
\qn{5} The dipole moment of a system of charge $+q$ distributed uniformly on an arc of radius $R$ subtending an angle $\pi/2$ at its centre where another charge $-q$ is placed is (In SI unit)\\
Given ( $q=1$C , $R=\sqrt8\,\pi$ m)
\qn{6} The two ends of a uniform thin rod of length $\sqrt2R$ and of mass $2\sqrt2$ kg can move without friction along a vertical circular path of radius $R$. The rod is released from the vertical position (ab). The force (in N) exerted by an end of the rod on the path when the rod passes the horizontal position (cd) is $10n$. Find `$n$'. ($g=10$ m/s$^2$)
\qfig{p2_phy_s3_q06_1}
\partsec{2}{CHEMISTRY}
\sect{Section-I}
\qn{1} One mole of an ideal monoatomic gas undergoes process from state At A to C as :
\[\underset{\substack{1\,\text{atm}\\300\text{K}}}{\text{State--A}}\xrightarrow[\text{process}]{\text{isochoric}}\underset{600\text{K}}{\text{State--B}}\xrightarrow[\text{Reversible Process}]{\text{Adiabatic}}\underset{64\,\text{atm}}{\text{State--C}}\]
Choose the correct option(s)
\begin{qoptions}
\opt{A}{Pressure of gas at state B is $2$ atm}
\opt{B}{$\Delta H_{AB}=900$ cal}
\opt{C}{$\Delta U_{AC}=6300$ cal}
\opt{D}{$W_{BC}=5400$ cal}
\end{qoptions}
\qn{2} Which of the following statements are correct ?
\begin{qoptions}
\opt{A}{At ordinary temperature oxidising strength of HClO $>$ HClO$_3$}
\opt{B}{Fluorine can decomposes water with liberation of ozonised oxygen.}
\opt{C}{HI is a stronger acid than HBr because of the low dissociation energy of HI.}
\opt{D}{In aqueous solution chlorine is a stronger oxidising agent than fluorine.}
\end{qoptions}
\qn{3} Which of the following carboxylic acid could be resolved by reaction with an enantiometrically pure chiral amine ?
\begin{qoptions}
\opt{A}{\qfigi{p2_chem_s1_q03_1}}
\opt{B}{\qfigi{p2_chem_s1_q03_B}}
\opt{C}{\qfigi{p2_chem_s1_q03_C}}
\opt{D}{\qfigi{p2_chem_s1_q03_D}}
\end{qoptions}
\qn{4} $0.1$ M H$_2$A Solution is titrated with $0.1$ M NaOH solution in the presence of methyl orange and phenolphthalein indicator, reaction as followed.
\[\mathrm{H_2A+NaOH\longrightarrow NaHA+H_2O}\ \text{(methyl orange)}\]
\[\mathrm{NaHA+NaOH\longrightarrow Na_2A+H_2O}\ \text{(Phenolphthalein)}\]
given that $K_{a1}(\mathrm{H_2A})=10^{-3}$, $K_{a2}=10^{-6}$\\
Select the correct analytical statement(s) about above titration.
\begin{qoptions}
\opt{A}{$100$ ml H$_2$A solution is titrated with $100$ ml NaOH in the presence of methyl orange, final pH $=4.5$}
\opt{B}{$100$ ml H$_2$A solution is titrated with $200$ ml NaOH in the presence of phenolphthalein, final pH $=9.26$}
\opt{C}{$100$ ml H$_2$A solution is titrated with $400$ ml NaOH in the presence of phenolphthalein , final pH $=12.6$}
\opt{D}{When $100$ ml H$_2$A solution is added with $50$ ml NaOH solution, buffer solution will be formed}
\end{qoptions}
\qn{5} Which of the following statements are correct :-
\begin{qoptions}
\opt{A}{Anhydrous MgCl$_2$ can not be prepared by heating MgCl$_2$.6H$_2$O}
\opt{B}{Anhydrous CaCl$_2$ is used in drying gases and organic compounds but not NH$_3$ or ethyl alcohol.}
\opt{C}{Solvay process is used for preparation of Na$_2$CO$_3$ but not for K$_2$CO$_3$}
\opt{D}{KOH is preferably used compared to NaOH for absorption of CO$_2$ because KHCO$_3$ formed is soluble where as NaHCO$_3$ is sparingly soluble and therefore choke the tubes of apparants used.}
\end{qoptions}
\qn{6} The product(s) obtained during reaction is/are
\qfig{p2_chem_s1_q06_1}
\begin{qoptions}
\opt{A}{\qfigi{p2_chem_s1_q06_A}}
\opt{B}{\qfigi{p2_chem_s1_q06_B}}
\opt{C}{\qfigi{p2_chem_s1_q06_C}}
\opt{D}{\qfigi{p2_chem_s1_q06_D}}
\end{qoptions}
\sect{Section-II}
\qn{1} By a sample of ground state atomic hydrogen, UV light of energy $\dfrac{13.6\times48}{49}\ \dfrac{\text{eV}}{\text{quanta}}$ is absorbed. How many different wavelengths will be observed Balmer region of hydrogen spectrum?
\qn{2} The spin only magnetic moment value (in Bohr magneton units) of Cr(CO)$_6$ is
\qn{3} \begin{center}
(a)~\qfigi{p2_chem_s2_q02_1}\qquad (b)~\qfigi{p2_chem_s2_q02_2}
\end{center}
X $=$ Total number of alkenes which on reaction with dil. H$_2$SO$_4$ forms (a) as major product.\\
Y $=$ Total number of alkenes which on reaction with dil. H$_2$SO$_4$ forms (b) as major product.\\
Then, find (Y $-$ X)
\qn{4} A real gas (X) has molar volume $1.0003$ times greater than ideal gas in experimental conditions, at working pressure and temperature $1$ bar and $300$ K respectively. Boyle's temperature of gas was observed $200$ K. What is value of excluded volume in litre/mole of gas (Consider value of $R=0.08$ bar litre K$^{-1}$ mole$^{-1}$)
\qn{5} Find the number of oxy acids having peroxy linkage\\
H$_4$P$_2$O$_6$, H$_4$P$_2$O$_8$, H$_4$P$_2$O$_7$, H$_2$SO$_5$, H$_4$P$_2$O$_5$, H$_3$P$_3$O$_9$, H$_6$P$_4$O$_{13}$, H$_2$S$_2$O$_7$, H$_2$S$_2$O$_3$, H$_2$S$_2$O$_5$, HClO$_4$
\qn{6} \begin{center}
$\begin{array}{l}\mathrm{CH{=}O}\\ |\\ \mathrm{CHOH}\\ |\\ \mathrm{CHOH}\\ |\\ \mathrm{CHOH}\\ |\\ \mathrm{CH_2OH}\end{array}$\ $\xrightarrow[\substack{\text{3.}\Delta\\\text{4. Na(Hg)/H}_2\text{SO}_4}]{\substack{\text{1. HCN}\\\text{2. Ba(OH)}_2\text{/H}_2\text{SO}_4}}$\ P
\end{center}
Number of carbon atoms present in P
\sect{Section-III}
\qn{1} KI$_{\text{(aq.)}}$ and KIO$_{3\text{(aq.)}}$ mixed together in presence of HCl to form I$_2$ and KCl and liberated I$_2$ required $100$ mL of $0.1$ M Na$_2$S$_2$O$_3$ solution for complete reduction of I$_2$.\\
If $x$ $=$ millimoles of KI used in the reaction.\\
and $y$ $=$ millimoles of HCl used in the reaction.\\
then value of $(3x-2y)$ is :
\qn{2} Find the number of ligands which are stronger ligand as compound to ethylenediamine(en)\\
NO$_2^-$; H$_2$O; NO$_3^-$; F$^-$; C$_2$O$_4^{2-}$; NH$_3$; Cl$^-$; CN$^-$
\qn{3} $0.436$ gm of acetyl derivative of a polyhydric alcohol requires $0.33$ gm KOH for hydrolysis. Calculate number of hydroxy groups in alcohol. (Molecular weight of polyhydric alcohol=$92$)
\qn{4} A sample of U$^{238}$ (half life $=4.5\times10^9$ yr) ore is found to contain $23.8$ g of U$^{238}$ and $20.6$ g of Pb$^{206}$. If the age of the ore is $A\times10^9$, find the value of $2A$.
\qn{5} Salt of iron (X) $\xrightarrow{\text{Na}_2\text{S}_2\text{O}_3}$ Violet complex (Y) $\xrightarrow{\text{On Standing}}$ Green solution (Z)\\
Write sum of the oxidation number of Fe in (X), (Y) and (Z).
\qn{6} Total number of correct statements are :\\
(i) Benzaldehyde cannot be obtained by Friedel craft acylation\\
(ii) Nitrobenzene is used as a solvent during Friedel craft reaction.\\
(iii) Friedel craft acylation requires higher concentration of catalyst than Friedel craft alkylation.\\
(iv) Benzaldehyde never reduces Fehling solution.\\
(v) (NH$_4$)$_2$S is used for selective reduction of only one $-$ NO$_2$ group into $-$NH$_2$\\
(vi) In coupling reaction attacking electrophile is $\mathrm{R{-}C{\equiv}O^{\oplus}}$ and $\mathrm{Ph{-}\overset{\oplus}{N}{\equiv}N}$\\
(vii) Polyalkylation is a disadvantage during monoalkylation of aromatic compound by Friedel craft reaction.\\
(viii) Direct nitration of aniline gives yield of para-nitro aniline.
\partsec{3}{MATHEMATICS}
\sect{Section-I}
\qn{1} If `$n$' denotes number of planes equidistant from four non-coplanar points. Then `$n$' is a factor of
\begin{qoptions}
\opt{A}{$12$}
\opt{B}{$14$}
\opt{C}{$24$}
\opt{D}{$42$}
\end{qoptions}
\qn{2} The curve satisfying the differential equation $\left(2x^2y-2y^4\right)dx+\left(2x^3+3xy^3\right)dy=0$ and passing through $(1,1)$ is given by $2\ell n(xy)+\dfrac{y^m}{x^n}=1$ $(m,n\in N)$, then
\begin{qoptions}
\opt{A}{$m+n=5$}
\opt{B}{$m+n=4$}
\opt{C}{$m\times n=6$}
\opt{D}{$m\times n=4$}
\end{qoptions}
\qn{3} Let $f$ be a differentiable function on $R$ and satisfying the integral equation
\[\int_0^xf(t)\,dt+\int_0^xt\cdot f(x-t)\,dt=-1+e^{-x},\]
for all $x\in R$, then
\begin{qoptions}
\opt{A}{$f(2)=e^{-2}$}
\opt{B}{$f(0)+f'(0)=1$}
\opt{C}{$f'(0)=2$}
\opt{D}{$f'(0)=1$}
\end{qoptions}
\qn{4} If $f(x)$ and $g(x)$ are differentiable functions for $x_1\le x\le x_2$ such that $f(x_1)=4,g(x_1)=1,f(x_2)=6$ and $g(x_2)=2$, then for some $C\in(x_1,x_2)$, $f'(C)=\lambda g'(C)$ where $\lambda$ equals
\begin{qoptions}
\opt{A}{Number of divisors of any prime number}
\opt{B}{Number of ways of putting $5$ identical balls in three identical boxes. (empty box not allowed)}
\opt{C}{Number of values of `$x$' satisfying the equation $e^{\sin x}-x=0$, $x\in R$}
\opt{D}{Number of into functions $f:A\to B$ where set $A$ has four distinct elements, and set $B$ has two distinct elements.}
\end{qoptions}
\qn{5} $f(x)=\begin{cases}a+\sin(\sin x);&x\ge0\\ \ell n(\cos x)+bx;&x<0\end{cases}$ is differentiable at $x=0$, then -
\begin{qoptions}
\opt{A}{$a=1$}
\opt{B}{$b=0$}
\opt{C}{$a=0$}
\opt{D}{$b=1$}
\end{qoptions}
\qn{6} If ${}^nC_1-\left(1+\dfrac12\right){}^nC_2+\left(1+\dfrac12+\dfrac13\right){}^nC_3\dots+(-1)^{n-1}\left(1+\dfrac12+\dfrac13+\dots+\dfrac1n\right){}^nC_n=\dfrac\lambda n$ then $\lambda\ (\lambda\in N)$ is less than
\begin{qoptions}
\opt{A}{$1$}
\opt{B}{$2$}
\opt{C}{$3$}
\opt{D}{$4$}
\end{qoptions}
\sect{Section-II}
\qn{1} If $f(x)=\begin{vmatrix}\cos x&\cos x&\sec^2x+\cot x\cdot\operatorname{cosec}x\\ \cos^2x&\cos^2x&\operatorname{cosec}^2x\\ 1&\cos^2x&\cos^2x\end{vmatrix}$ then find the value of $\left|\displaystyle\int_0^{\pi/2}f(x)\,dx\right|$
\qn{2} Let $P$ be the product of the non-real roots of the equation $x^4-4x^3+6x^2-4x=2020$, then find $P$.
\qn{3} Let $a,b,c$ be positive real numbers, such that, $\dfrac1{1+a}+\dfrac1{1+b}+\dfrac1{1+c}\le1$, Now if the least value of $\left(1+a^2\right)\cdot\left(1+b^2\right)\cdot\left(1+c^2\right)$ is $\lambda$, then $\dfrac\lambda4$ equals.
\qn{4} If $z_1^3-3z_1z_2^2=2$ and $3z_2z_1^2-z_2^3=11$, then find the value of $\dfrac{\left|z_1^2+z_2^2\right|}{4}$
\qn{5} The slope of one of the common tangents to the hyperbolas $\dfrac{x^2}{16}-\dfrac{y^2}9=1$ and $\dfrac{y^2}{16}-\dfrac{x^2}9=1$ is $m$, then $|m|$ is
\qn{6} If $\phi_1,\phi_2,\phi_3,\phi_4$ are the eccentric angles of four concyclic points on the ellipse $\dfrac{x^2}{a^2}+\dfrac{y^2}{b^2}=1$, then find the value of $\dfrac74\cos\left(\phi_1+\phi_2+\phi_3+\phi_4\right)$.
\sect{Section-III}
\qn{1} How many tangents to the circle $x^2+y^2=3$ are there which are normal to the ellipse $\dfrac{x^2}{9}+\dfrac{y^2}{4}=1$
\qn{2} The lines $y=mx$ and $y=nx$ intersect the line $x+y=4$ at points A and B respectively such that OA, AB and OB are in Arithmetic progression, where `O' is the origin then $|m-n|\le K|1+mn|$. Find the least possible value of $[K]$ where $[\cdot]$ denotes greatest integer function.
\qn{3} Find the number of values of `$x$' satisfying $\sin x\cdot\sin\left(\dfrac1x\right)-1=0$, $x\in R-\{0\}$
\qn{4} Let an ordered pair A be defined as A $(x,y)$ where $x\in$ prime number, such that $x<10$ and $y\in$ natural numbers and $y\le10$. If the probability that the ordered pair A satisfies the relation $x^2-3y^2=1$ is P then $60$ P equals.
\qn{5} Let A and B are two non-singular matrices such that $A^6=I$ and $AB^2=BA$. If the least value of $m$ such that $B^m=I$ $(B\neq I)$ $(m\in N)$ is K then find the value of $\dfrac K9$
\qn{6} If $8\displaystyle\int_0^{\pi/2}\dfrac{x^2(\cos x-\sin x)\,dx}{(\cos x+\sin x)^3}=\pi(k-\pi)$, then $k$ is equal to
\end{multicols}
\end{document}
